% =============================================================================
%  Préambule du rapport final InfraReady (LuaLaTeX)
% =============================================================================

% --- Langue et polices -------------------------------------------------------
\usepackage[french]{babel}
\usepackage{fontspec}
\directlua{luaotfload.add_fallback("repli", {"DejaVuSans:mode=harf;", "DejaVuSansMono:mode=harf;"})}
\setmainfont{TeX Gyre Pagella}[RawFeature={fallback=repli}]
\setsansfont{TeX Gyre Heros}[RawFeature={fallback=repli}]
\setmonofont{DejaVu Sans Mono}[Scale=0.84]
\usepackage[final]{microtype}

% --- Mise en page --------------------------------------------------------------
\usepackage[a4paper,top=2.4cm,bottom=2.6cm,left=2.3cm,right=2.3cm,headheight=15pt]{geometry}
\usepackage[autooneside=false,automark]{scrlayer-scrpage}
\clearpairofpagestyles
\ihead{\small\sffamily InfraReady — Rapport final}
\ohead{\small\sffamily\leftmark}
\cfoot*{\small\sffamily\pagemark}
\setkomafont{pageheadfoot}{\color{gris}}
\renewcommand*{\chapterpagestyle}{plain.scrheadings}

% --- Couleurs ------------------------------------------------------------------
\usepackage[table]{xcolor}
\definecolor{bleu}{HTML}{1F4E79}
\definecolor{bleuclair}{HTML}{DCE8F5}
\definecolor{vert}{HTML}{1A7F37}
\definecolor{vertclair}{HTML}{DDF4E4}
\definecolor{orange}{HTML}{9A6700}
\definecolor{orangeclair}{HTML}{FFF3D6}
\definecolor{rouge}{HTML}{B42318}
\definecolor{rougeclair}{HTML}{FDE7E5}
\definecolor{gris}{HTML}{57606A}
\definecolor{grisclair}{HTML}{F3F4F6}

\setkomafont{disposition}{\sffamily\bfseries\color{bleu}}
\setkomafont{chapter}{\Huge}
\setkomafont{descriptionlabel}{\sffamily\bfseries}
\RedeclareSectionCommand[beforeskip=-1.2\baselineskip,afterskip=.6\baselineskip]{section}
\RedeclareSectionCommand[beforeskip=-1\baselineskip,afterskip=.4\baselineskip]{subsection}

% --- Tableaux, figures, listes -------------------------------------------------
\usepackage{array,booktabs,tabularx,longtable,colortbl}
\usepackage{graphicx}
\graphicspath{{rapport/figures/}{docs/preuves/}}
\usepackage[font=small,labelfont={sf,bf},format=hang,justification=raggedright,singlelinecheck=false]{caption}
\usepackage{float}
\usepackage{enumitem}
\setlist{itemsep=2pt,topsep=3pt}
\usepackage{pdflscape}
\usepackage{tikz}
\usetikzlibrary{positioning,shapes.geometric}
\renewcommand{\arraystretch}{1.25}
\newcolumntype{L}{>{\raggedright\arraybackslash}X}
\newcolumntype{P}[1]{>{\raggedright\arraybackslash}p{#1}}

% --- Hyperliens -----------------------------------------------------------------
\usepackage[hidelinks,pdfusetitle,bookmarksnumbered]{hyperref}
\hypersetup{
  pdftitle={InfraReady — Rapport final},
  pdfauthor={MIR Mahan, Jolan Allen, Kenzi Bendjelloul},
  pdfsubject={Projet fil rouge B3 Cybersécurité : standardisation, automatisation et supervision},
  pdfkeywords={Proxmox, Active Directory, Ansible, Prometheus, Grafana, GitHub Actions}
}
\usepackage{bookmark}

% --- Commandes du rapport -------------------------------------------------------
% Code en ligne : le contenu est pris tel quel (soulignés, points, barres obliques…) et peut être
% coupé en fin de ligne après _ / . - (rapport/code-secable.lua), utile dans les colonnes étroites.
\directlua{dofile(kpse.find_file('rapport/code-secable.lua') or 'rapport/code-secable.lua')}
\DeclareRobustCommand{\code}[1]{\texttt{\directlua{codesecable("\luaescapestring{\detokenize{#1}}")}}}
% Renvoi à une preuve du dépôt (docs/preuves/)
\DeclareRobustCommand{\preuve}[1]{\textsf{\textbf{\color{bleu}#1}}}
% Statuts
\newcommand{\ok}{\textsf{\textbf{\color{vert}OK}}}
\newcommand{\fait}{\textsf{\textbf{\color{vert}Fait}}}
\newcommand{\partiel}{\textsf{\textbf{\color{orange}Partiel}}}
\newcommand{\afaire}{\textsf{\textbf{\color{rouge}À faire}}}

% Encadré coloré (non sécable) : \encadre{couleur fond}{couleur trait}{titre}{texte}
\newcommand{\encadre}[4]{%
  \par\smallskip\noindent
  \begin{tikzpicture}
    \node[fill=#1, draw=#2, line width=0.6pt, rounded corners=3pt,
          inner sep=7pt, text width=\linewidth-26pt, align=left] (b) {%
      {\sffamily\bfseries\color{#2}#3}\par\smallskip #4};
    \fill[#2] ([xshift=0pt]b.north west) rectangle ([xshift=3pt]b.south west);
  \end{tikzpicture}\par\smallskip}
\newcommand{\info}[2]{\encadre{bleuclair}{bleu}{#1}{#2}}
\newcommand{\succes}[2]{\encadre{vertclair}{vert}{#1}{#2}}
\newcommand{\attention}[2]{\encadre{orangeclair}{orange}{#1}{#2}}
\newcommand{\alerte}[2]{\encadre{rougeclair}{rouge}{#1}{#2}}

% Figure standard : \figure{fichier}{largeur relative}{légende}{étiquette}
\newcommand{\illustration}[4]{%
  \begin{figure}[H]
    \centering
    \includegraphics[width=#2\linewidth,height=0.72\textheight,keepaspectratio]{#1}
    \caption{#3}\label{#4}
  \end{figure}}

% Fiche de phase (objectif, responsables, dates, statut)
\newcommand{\fiche}[4]{%
  \info{Objectif}{#1\par\smallskip
  \small\textbf{Responsables :} #2 \quad·\quad \textbf{Dates :} #3 \quad·\quad \textbf{Statut :} #4}}
